\documentclass[11pt]{article}

\usepackage[T1]{fontenc}
\usepackage{lmodern}
\usepackage{amsmath,amssymb,amsthm}
\usepackage[margin=1in]{geometry}
\usepackage{microtype}
\usepackage[hidelinks]{hyperref}
\hypersetup{
  pdftitle={A bound for the number of maximum modulus points},
  pdfauthor={Qiyuan Gu}
}

\newtheorem{theorem}{Theorem}[section]
\newtheorem{lemma}[theorem]{Lemma}
\newtheorem{proposition}[theorem]{Proposition}
\theoremstyle{remark}
\newtheorem{remark}[theorem]{Remark}
\numberwithin{equation}{section}
\DeclareMathOperator{\ord}{ord}
\newcommand{\C}{\mathbb C}
\newcommand{\R}{\mathbb R}
\newcommand{\X}{\mathcal X}
\newcommand{\B}{\mathcal B}

\title{A bound for the number of maximum modulus points}
\author{Qiyuan Gu\\
University of Chicago\\
\href{mailto:phoenix1203@uchicago.edu}{\texttt{phoenix1203@uchicago.edu}}}
\date{}

\begin{document}
\maketitle

\begin{abstract}
For a non-monomial entire function \(f\), let \(\nu_f(r)\) be the
number of points on \(|z|=r\) at which \(|f|\) attains its maximum.
Write \(f(z)=z^m(c_0+c_kz^k+O(z^{k+1}))\), where
\(c_0c_k\ne0\) and \(k\ge1\).
We prove that \(\nu_f(r)\le2k\) outside a countable set of radii.
In particular, \(\nu_f(r)\) cannot tend to infinity, answering the
second question of Erd\H{o}s on maximum modulus points.
The proof studies the correspondence
\(A(z)=\overline{A(\overline w)}\), where \(A=zf'/f\), using
implicit-function continuation, growth in a direct tract, and a finite
proper map of analytic curves.
\end{abstract}

\section{Introduction}

For a nonzero entire function \(f\), put
\[
 M(r,f)=\max_{|z|=r}|f(z)|,
 \qquad
 \nu_f(r)=\#\{z:|z|=r,\ |f(z)|=M(r,f)\}.
\]
If \(f\) is not a monomial, then \(\nu_f(r)\) is finite for every
\(r>0\). Erd\H{o}s asked whether a single such function can satisfy
\[
 \limsup_{r\to\infty}\nu_f(r)=\infty,
 \qquad\text{and whether it can satisfy}\qquad
 \liminf_{r\to\infty}\nu_f(r)=\infty;
\]
see \cite{Erdos1117,GP} for the problem and its history.
Herzog and Piranian \cite{HP} answered the first question affirmatively.
Pardo-Sim\'on and Sixsmith \cite{PS} obtained such an example of finite
order. Gl\"ucksam and Pardo-Sim\'on \cite{GP} constructed an approximate
analogue of the second property, with many separated arcs on which
the modulus is close to its maximum.

There is also a classical local theory near the origin. Hayman
\cite{Hayman} proved that for \(f\) as in \eqref{eq:taylor}, the
maximum modulus set near zero consists of at most \(k\) analytic
curves. Evdoridou, Pardo-Sim\'on, and Sixsmith \cite{EPS} refined
this description, determining the exact number of local maximum
curves outside an algebraically defined exceptional class and showing
that each such curve contains exactly one point of each sufficiently
small positive modulus. In particular, \(\nu_f(r)\le k\) for all
sufficiently small \(r\).

We prove that the answer to the second question is negative.

\begin{theorem}\label{thm:main}
Let \(f\) be a nonzero entire function that is not a monomial, and write
\begin{equation}\label{eq:taylor}
 f(z)=z^m\bigl(c_0+c_kz^k+O(z^{k+1})\bigr),
 \qquad m\ge0,\quad k\ge1,\quad c_0c_k\ne0.
\end{equation}
There is a countable set \(E\subset(0,\infty)\) such that
\begin{equation}\label{eq:main-bound}
 \nu_f(r)\le2k\qquad(r\notin E).
\end{equation}
Consequently,
\[
 \liminf_{r\to\infty}\nu_f(r)\le2k<\infty.
\]
\end{theorem}

The integer \(k\) is the difference between the degrees of the first
two nonzero terms of \(f\). The factor \(2\) is attained by an explicit
example given in Remark~\ref{rem:sharp}.


The countable-exception bound allows the unbounded upper limits constructed
in \cite{HP,PS}: their large counts may occur on the exceptional radii.
The separated arcs in \cite{GP} concern values close to the maximum, so they
need not produce additional points attaining it. The local bound \(k\)
near the origin from \cite{Hayman,EPS} is also consistent with the bound
\(2k\) for nonexceptional radii here.

Set \(A=zf'/f\). At radii where \(\log M(r,f)\) is differentiable
with respect to \(\log r\), all maximum modulus points have the same
real value of \(A\). Hence they determine points \((z,\overline z)\)
on the correspondence \(A(z)=\overline{A(\overline w)}\) with the same
values of \(zw\) and \(A(z)\). We show that every irreducible component
of this correspondence contains points with \(|zw|\) arbitrarily small.
Since \(A-m\) has local degree \(k\) at the origin, the generic fibres
of the resulting finite proper map have size at most \(2k\).

\section{Analytic preliminaries}\label{sec:facts}

\begin{lemma}[Sto\"ilow; see \cite{Eremenko}]\label{lem:continuation}
Let \(F\) be entire on \(\C^2\), and suppose that
\[
 F(a,b)=0,\qquad F_w(a,b)\ne0.
\]
The projection to the \(z\)-plane of the maximal analytic continuation
of the implicit germ \(w=w(z)\) through \((a,b)\) is dense in \(\C\).
\end{lemma}

Following Bergweiler, Rippon, and Stallard
\cite[Definition 2.1]{BRS}, an unbounded domain \(D\subset\C\), with
unbounded complement and piecewise smooth boundary, is a
\emph{direct tract} of \(Q\) if \(Q\) is holomorphic in \(D\),
continuous on its closure, and
\[
 |Q|=a\quad\text{on }\partial D,
 \qquad |Q|>a\quad\text{in }D
\]
for some \(a>0\).

\begin{lemma}[Fuchs; see {\cite[Theorem 2.1]{BRS}}]
\label{lem:tract}
Let \(D\) be a direct tract of \(Q\), with boundary value \(a\), and put
\[
 u(w)=
 \begin{cases}
  \log\bigl(|Q(w)|/a\bigr),&w\in D,\\
  0,&w\notin D.
 \end{cases}
\]
Then \(u\) is a continuous subharmonic function on \(\C\), and
\begin{equation}\label{eq:tract-growth}
 \frac{\max_{|w|=R}u(w)}{\log R}\longrightarrow\infty
 \qquad(R\to\infty).
\end{equation}
\end{lemma}

\section{Small products on the correspondence}\label{sec:products}

Fix \(f,m,k\) as in Theorem~\ref{thm:main}. Write \(f=z^m h\), where
\(h(0)\ne0\), and define the meromorphic functions
\begin{equation}\label{eq:AB}
 A(z)=\frac{zf'(z)}{f(z)}=m+\frac{zh'(z)}{h(z)},
 \qquad B(w)=A^\#(w):=\overline{A(\overline w)}.
\end{equation}
The removable singularity of \(A\) at zero is filled in.
The function \(A\) is nonconstant, since otherwise \(f\) would be a
monomial. Moreover, \eqref{eq:taylor} gives
\[
 A(z)-m=\frac{kc_k}{c_0}z^k+O(z^{k+1}).
\]
In particular, \(k=\ord_0(A-m)\).

There is a holomorphic local coordinate \(\chi\) at zero such that
\begin{equation}\label{eq:coordinate}
 \chi(0)=0,\qquad\chi'(0)\ne0,
 \qquad A(z)-m=\chi(z)^k.
\end{equation}
Choose \(\delta\in(0,1)\) small enough that
\[
 V=\chi^{-1}\{\xi:|\xi|<\delta^{1/k}\}
\]
has compact closure in the coordinate neighbourhood.
For every \(|\lambda-m|<\delta\), the equation \(A(z)=\lambda\)
has exactly \(k\) roots in \(V\), counted with multiplicity.
They are distinct when \(\lambda\ne m\), and there is a constant
\(C\) such that all of them satisfy
\begin{equation}\label{eq:root-bound}
 |z|\le C|\lambda-m|^{1/k}.
\end{equation}
The corresponding statements hold for \(B\) on
\(V^\#=\{\overline z:z\in V\}\).

Choose \(\delta\) also so that it is not the modulus of any finite,
nonzero critical value of \(B-m\). Then
\(\{|B-m|=\delta\}\) is a locally finite union of smooth curves.

Choose entire functions \(p,q\) with no common zero such that
\(A=p/q\) and \(q(0)\ne0\). Set
\begin{equation}\label{eq:correspondence}
 F(z,w)=p(z)q^\#(w)-p^\#(w)q(z),
 \qquad \X=\{(z,w)\in\C^2:F(z,w)=0\}.
\end{equation}
All analytic sets in this paper are given their reduced structure.
Away from poles, the equation in \eqref{eq:correspondence} is
\(A(z)=B(w)\).
Since \(p,q\) have no common zero and \(A\) is nonconstant,
\(\X\) has no horizontal or vertical irreducible component.
At points with finite common value,
\[
 F_w(z,w)=-q(z)q^\#(w)B'(w).
\]
It follows that \(F_w\ne0\) at a general point of each irreducible
component: the coordinate \(w\) is nonconstant there, whereas the
zeros of \(B'\) are discrete.

\begin{lemma}\label{lem:small-products}
Every irreducible component \(X\) of \(\X\) contains a sequence
\((z_j,w_j)\) such that
\begin{equation}\label{eq:small-products}
 z_jw_j\ne0,\qquad
 A(z_j)=B(w_j)\in\C\setminus\{m\},
 \qquad |z_jw_j|\longrightarrow0.
\end{equation}
\end{lemma}

\begin{proof}
Choose a regular point of \(X\) at which \(F_w\ne0\) and apply
Lemma~\ref{lem:continuation} to the corresponding implicit germ.
Since its projection is dense, the continuation reaches a point
\((z_0,w_0)\in X\) with
\[
 z_0\in V\setminus\{0\},\qquad
 0<|A(z_0)-m|<\delta.
\]
Let \(U\) be the connected component of \(\{|B-m|<\delta\}\)
containing \(w_0\).
In \(V\times U\), consider the analytic curve
\begin{equation}\label{eq:root-relation}
 \chi(z)^k=B(w)-m.
\end{equation}
Its projection to \(U\) is finite and proper. Indeed, over a compact
subset of \(U\), the right-hand side remains in a compact subset of
\(\{\xi:|\xi|<\delta\}\), so all its \(k\)-th roots remain in a
compact subset of the coordinate disc.

Since \(A'(z_0)\ne0\), the curve \eqref{eq:root-relation} is smooth
at \((z_0,w_0)\). Let \(Y\) be its irreducible component through this
point. Near \((z_0,w_0)\), it is the same branch as \(X\), hence
\(Y\subset X\) by the identity theorem for analytic sets. The
projection \(Y\to U\) is finite, proper, and nonconstant. Its image is
therefore both open and closed in \(U\), so the projection is
surjective. Thus every \(w\in U\) has a lift \((z,w)\in Y\), with
\(z\) satisfying \eqref{eq:root-bound}.

If \(B(b)=m\) for some \(b\in U\), choose \(w_j\to b\) in \(U\),
with \(w_j\ne0\) and \(B(w_j)\ne m\), and lift these points to \(Y\).
Equation~\eqref{eq:root-bound} gives \(z_j\to0\), and hence
\(z_jw_j\to0\). This proves the lemma in this case.

Suppose now that \(B-m\) has no zero in \(U\), and put
\begin{equation}\label{eq:Q}
 Q(w)=\frac{\delta}{B(w)-m}.
\end{equation}
Every finite boundary point of \(U\) satisfies \(|B-m|=\delta\).
Thus \(Q\) is holomorphic in \(U\), continuous on \(\overline U\),
and satisfies
\[
 |Q|>1\quad\text{in }U,
 \qquad |Q|=1\quad\text{on }\partial U.
\]
The maximum principle shows that \(U\) is unbounded.

We claim that its complement is also unbounded.
If \(f\) has a zero other than the origin, then for every sufficiently
large \(R\) whose circle avoids the zeros of \(f\), the argument
principle gives
\begin{equation}\label{eq:mean}
 \frac1{2\pi}\int_0^{2\pi}
 \bigl(B(Re^{i\theta})-m\bigr)\,d\theta
 =n(R,f)-m\ge1,
\end{equation}
where \(n(R,f)\) counts zeros in \(|z|<R\), with multiplicity.
Here \(B(w)=w(f^\#)'(w)/f^\#(w)\), and \(f\) and \(f^\#\) have
the same zero counts on centred discs.
Since \(\delta<1\), such a circle cannot be contained in
\(\{|B-m|<\delta\}\). A circle containing a pole of \(B\) is not
contained in that set either.
If \(f\) has no zero other than possibly the origin, then \(B-m\)
is a nonconstant entire function. Its maximum modulus eventually
exceeds \(\delta\), giving the same conclusion for all large circles.
In either case, \(\C\setminus U\) is unbounded.

Consequently, \(U\) is a direct tract of \(Q\) with boundary value
one. Define \(u=\log|Q|\) in \(U\), and \(u=0\) outside \(U\).
For each sufficiently large \(R\), choose \(w_R\) on \(|w|=R\)
at which \(u\) attains its maximum.
By Lemma~\ref{lem:tract}, this maximum is positive and
\[
 \frac{u(w_R)}{\log R}\longrightarrow\infty.
\]
In particular, \(w_R\in U\).
Choose a lift \((z_R,w_R)\in Y\). Equations
\eqref{eq:root-bound} and \eqref{eq:Q} yield
\begin{equation}\label{eq:product-decay}
 |z_Rw_R|
 \le C\delta^{1/k}R\exp\!\left(-\frac{u(w_R)}k\right)
 \longrightarrow0.
\end{equation}
Since \(B-m\) has no zero in \(U\), these lifts have nonzero
coordinates and common value different from \(m\).
They satisfy \eqref{eq:small-products}.
\end{proof}

\section{A finite map of analytic curves}\label{sec:finite}

Remove the zero coordinates, poles, and common value \(m\), and set
\begin{align*}
 \X^*&=\{(z,w)\in\X:zw\ne0,\ A(z)=B(w)\in\C\setminus\{m\}\},\\
 \B&=\C^*\times(\C\setminus\{m\}).
\end{align*}
On each irreducible component of \(\X\), the deleted locus is
discrete. Thus every irreducible component of \(\X^*\) is obtained
from one of \(\X\) by deleting a discrete set.
Consider the holomorphic map
\begin{equation}\label{eq:Phi}
 \Phi:\X^*\longrightarrow\B,
 \qquad \Phi(z,w)=(zw,A(z))=(s,\lambda).
\end{equation}

\begin{lemma}\label{lem:proper}
The map \(\Phi\) is proper and has finite fibres.
\end{lemma}

\begin{proof}
Suppose that \((s,\lambda)\) remains in a compact subset of \(\B\).
Then \(|s|\) is bounded above and away from zero, while \(\lambda\)
is bounded and stays a positive distance from \(m\).
If a sequence of preimages had \(z\to\infty\), then
\(w=s/z\to0\), and therefore \(\lambda=B(w)\to m\), a contradiction.
If \(z\to0\), then \(\lambda=A(z)\to m\), again a contradiction.
The same arguments apply with \(z\) and \(w\) interchanged.
Neither coordinate can approach a pole, since \(\lambda\) is bounded.
It follows that the inverse image of the compact set is compact in
\(\X^*\).

For fixed \((s,\lambda)\), the possible \(z\)-coordinates form a
compact subset of the discrete zero set of \(A-\lambda\).
There are therefore only finitely many of them, and each determines
\(w=s/z\) uniquely.
\end{proof}

\begin{proposition}\label{prop:degree}
The image \(\Gamma=\Phi(\X^*)\) is a closed analytic curve in \(\B\).
There is a locally finite set \(T\subset\Gamma\) such that
\begin{equation}\label{eq:fibre-bound}
 \#\Phi^{-1}(s,\lambda)\le2k
 \qquad\bigl((s,\lambda)\in\Gamma\setminus T\bigr).
\end{equation}
\end{proposition}

\begin{proof}
By Remmert's proper mapping theorem and Lemma~\ref{lem:proper},
\(\Gamma\) is a closed analytic curve.
Let \(T\) consist of the singular points of \(\Gamma\) together with
the images of the singular and ramification points of \(\X^*\), where
ramification is understood on the normalization of each source
component. The map is nonconstant on every source component, so its
ramification locus is discrete; properness then makes these exceptional
sets locally finite in \(\Gamma\).

For each irreducible component \(\Gamma_i\) of \(\Gamma\), the map
\[
 \Phi^{-1}(\Gamma_i\setminus T)\longrightarrow\Gamma_i\setminus T
\]
is a finite covering. The normalization of \(\Gamma_i\) is a connected
Riemann surface, and deleting a locally finite set leaves it connected.
Since \(T\) contains the singular points of \(\Gamma\), it follows that
\(\Gamma_i\setminus T\) is connected. Hence there is an integer
\(d_i\) such that
\begin{equation}\label{eq:degree}
 \#\Phi^{-1}(s,\lambda)=d_i
 \qquad\bigl((s,\lambda)\in\Gamma_i\setminus T\bigr).
\end{equation}

Choose \(\eta>0\) so that the disc \(\{|z|<\eta\}\) is contained
in both \(V\) and \(V^\#\).
If \(0<|s|<\eta^2\), every preimage satisfies
\[
 \min\{|z|,|w|\}<\eta.
\]
There are at most \(k\) preimages with \(|z|<\eta\): the existence
of such a coordinate implies \(|\lambda-m|<\delta\), and the local
root count in Section~\ref{sec:products} applies. Each root determines
the other coordinate from \(w=s/z\).
Likewise there are at most \(k\) preimages with \(|w|<\eta\).
Hence
\begin{equation}\label{eq:small-fibre}
 \#\Phi^{-1}(s,\lambda)\le2k
 \qquad(0<|s|<\eta^2).
\end{equation}

The image of each irreducible component of \(\X^*\) is an
irreducible component of \(\Gamma\), and every component of
\(\Gamma\) arises in this way.
Lemma~\ref{lem:small-products} therefore shows that each \(\Gamma_i\)
meets \(\{0<|s|<\eta^2\}\).
This intersection is a nonempty open subset of a curve, so it is
not contained in \(T\).
Evaluating \eqref{eq:degree} at a point of this intersection and
using \eqref{eq:small-fibre} gives \(d_i\le2k\).
Equation~\eqref{eq:fibre-bound} follows for every component.
\end{proof}

\section{Maximum modulus points}\label{sec:maxima}

\begin{proof}[Proof of Theorem~\ref{thm:main}]
Set
\[
 v(x)=\log M(e^x,f),\qquad x\in\R.
\]
Hadamard's three-circles theorem makes \(v\) convex.
Its set \(N\) of nondifferentiability points is therefore at most
countable.
Suppose that \(x\notin N\), let \(r=e^x\), and let
\(z=re^{i\theta}\) be a maximum modulus point.
Since \(f(z)\ne0\), angular differentiation gives
\begin{equation}\label{eq:angular}
 0=\frac{\partial}{\partial\theta}\log|f(re^{i\theta})|
 =-\operatorname{Im}A(z).
\end{equation}
For this fixed \(\theta\), the function
\(t\mapsto\log|f(e^{t+i\theta})|\) lies below \(v(t)\) and agrees
with it at \(t=x\). Differentiating at this contact point yields
\begin{equation}\label{eq:radial}
 \operatorname{Re}A(z)=v'(x).
\end{equation}
Thus all maximum modulus points on the circle have the same real
value \(A(z)=\lambda:=v'(x)\).

In fact \(\lambda>m\).
To see this, write
\[
 b(x)=v(x)-mx=\log M(e^x,h).
\]
The function \(b\) is convex and strictly increasing, since the maximum
modulus of a nonconstant entire function strictly increases with radius.
For any \(t<x\), the convex secant-slope inequality gives
\[
 0<\frac{b(x)-b(t)}{x-t}\le b'(x).
\]
Hence \(\lambda=m+b'(x)>m\).

For every maximum modulus point \(z\) on \(|z|=r\), we now have
\begin{equation}\label{eq:maximal-fibre}
 (z,\overline z)\in\X^*,
 \qquad\Phi(z,\overline z)=(r^2,\lambda).
\end{equation}
Distinct maximum points give distinct points of this fibre.
If \((r^2,\lambda)\notin T\), Proposition~\ref{prop:degree} gives
\(\nu_f(r)\le2k\).

Since \(T\) is locally finite, it is countable. Define
\[
 E=\{e^x:x\in N\}\ \cup\
   \{\sqrt{s}:s>0\text{ and }(s,\lambda)\in T
      \text{ for some }\lambda\}.
\]
This is a countable set of positive radii, and
\eqref{eq:main-bound} holds for every \(r\notin E\).
Since every interval \((R,\infty)\) contains radii outside \(E\),
the asserted bound on the lower limit follows.
\end{proof}

\begin{remark}\label{rem:sharp}
For \(k\ge1\), the function
\(f(z)=\exp(2z^k-z^{2k})\) has first two nonzero Taylor terms
\(1+2z^k\). Put \(R=r^k\) and \(\phi=k\theta\). Then
\[
 \log|f(re^{i\theta})|
 =R^2+\frac12
   -2R^2\left(\cos\phi-\frac1{2R}\right)^2.
\]
For \(0<R\le1/2\), the maximum is attained when
\(\cos(k\theta)=1\), giving \(k\) distinct points. For \(R>1/2\),
it is attained precisely when \(\cos(k\theta)=1/(2R)\), giving
\(2k\) distinct points. Thus the bound in Theorem~\ref{thm:main}
is sharp for every \(k\).
\end{remark}

\section*{Use of generative AI}

GPT-6 Astra was used to generate the mathematical proofs and draft the
manuscript. GPT-5.6 Sol and Claude Opus 5 were used only for editorial
review of the exposition and did not contribute to the mathematical
arguments. The author reviewed the final manuscript and takes full
responsibility for its content.

\begin{thebibliography}{9}

\bibitem{BRS}
W.~Bergweiler, P.~J. Rippon, and G.~M. Stallard,
Dynamics of meromorphic functions with direct or logarithmic
singularities,
\emph{Proceedings of the London Mathematical Society} (3)
\textbf{97} (2008), 368--400.
\href{https://arxiv.org/abs/0704.2712}{arXiv:0704.2712}.

\bibitem{Erdos1117}
T.~F. Bloom,
\emph{Erd\H{o}s Problem \#1117},
\url{https://www.erdosproblems.com/1117}.
Accessed 5 September 2026.

\bibitem{Eremenko}
A.~Eremenko,
\emph{Singularities of implicit functions},
note, 4 April 2015.
\url{https://www.math.purdue.edu/~eremenko/dvi/gross2.pdf}.

\bibitem{EPS}
V.~Evdoridou, L.~Pardo-Sim\'on, and D.~J. Sixsmith,
On a result of Hayman concerning the maximum modulus set,
\emph{Computational Methods and Function Theory} \textbf{21} (2021),
779--795.
\href{https://arxiv.org/abs/2012.07409}{arXiv:2012.07409}.

\bibitem{GP}
A.~Gl\"ucksam and L.~Pardo-Sim\'on,
An approximate solution to Erdos' maximum modulus points problem,
\emph{Journal of Mathematical Analysis and Applications}
\textbf{531} (2024), no.~1, Article 127768.
\href{https://arxiv.org/abs/2208.11154}{arXiv:2208.11154}.

\bibitem{Hayman}
W.~K. Hayman,
A characterization of the maximum modulus of functions regular at the origin,
\emph{Journal d'Analyse Math\'ematique} \textbf{1} (1951), 135--154.

\bibitem{HP}
F.~Herzog and G.~Piranian,
The counting function for points of maximum modulus,
in \emph{Entire Functions and Related Parts of Analysis},
Proceedings of Symposia in Pure Mathematics, vol.~11,
American Mathematical Society, Providence, RI, 1968, pp.~240--243.

\bibitem{PS}
L.~Pardo-Sim\'on and D.~J. Sixsmith,
\emph{A finite-order answer to a problem of Erd\H{o}s on maximum
modulus points},
preprint, 2026.
\href{https://arxiv.org/abs/2607.09462}{arXiv:2607.09462}.

\end{thebibliography}

\end{document}
